\textbf{Phần 1. Trắc nghiệm lựa chọn}
\setcounterpageref{ex}{0}

\begin{ex} %[2D2N11]
Kết quả của phép tính $11 - 2$ là
\choice
{$8$}
{$10$}
{$7$}
{\True $9$}
\loigiai{
$11 - 2 = 9$.\\
\textbf{Đáp án của câu hỏi là D}
}
\end{ex}

\begin{ex} %[2D2N11]
Kết quả của phép tính $12 - 3$ là
\choice
{$8$}
{\True $9$}
{$10$}
{$7$}
\loigiai{
$12 - 3 = 9$.\\
\textbf{Đáp án của câu hỏi là B}
}
\end{ex}

\begin{ex} %[2D2H11]
Kết quả của phép tính $11 - 5$ là
\choice
{\True $6$}
{$5$}
{$7$}
{$8$}
\loigiai{
$11 - 5 = 6$.\\
\textbf{Đáp án của câu hỏi là A}
}
\end{ex}

\begin{ex} %[2D2H11]
Kết quả của phép tính $13 - 7$ là
\choice
{$5$}
{$7$}
{$8$}
{\True $6$}
\loigiai{
$13 - 7 = 6$.\\
\textbf{Đáp án của câu hỏi là D}
}
\end{ex}

\begin{ex} %[2D2H11]
Kết quả của phép tính $14 - 8$ là
\choice
{\True $6$}
{$5$}
{$7$}
{$8$}
\loigiai{
$14 - 8 = 6$.\\
\textbf{Đáp án của câu hỏi là A}
}
\end{ex}

\begin{ex} %[2D2H11]
Kết quả của phép tính $15 - 9$ là
\choice
{$5$}
{$7$}
{\True $6$}
{$8$}
\loigiai{
$15 - 9 = 6$.\\
\textbf{Đáp án của câu hỏi là C}
}
\end{ex}

\begin{ex} %[2D2V11]
Điền dấu thích hợp: $13 - 6 \dots 14 - 8$.
\choice
{$<$}
{$=$}
{Không so sánh được}
{\True $>$}
\loigiai{
$13 - 6 = 7$ và $14 - 8 = 6$. Vì $7 > 6$ nên điền dấu $>$.\\
\textbf{Đáp án của câu hỏi là D}
}
\end{ex}

\begin{ex} %[2D2V11]
Tìm $x$ biết: $x + 6 = 15$.
\choice
{$8$}
{\True $9$}
{$10$}
{$7$}
\loigiai{
$x = 15 - 6 = 9$.\\
\textbf{Đáp án của câu hỏi là B}
}
\end{ex}

\begin{ex} %[2D2V11]
Có 16 quả táo, ăn hết 9 quả táo. Còn lại số quả táo là
\choice
{$6$ quả}
{$8$ quả}
{\True $7$ quả}
{$9$ quả}
\loigiai{
$16 - 9 = 7$ quả.\\
\textbf{Đáp án của câu hỏi là C}
}
\end{ex}

\begin{ex} %[2D2V11]
Tính giá trị biểu thức: $17 - 8 + 3$.
\choice
{\True $12$}
{$11$}
{$13$}
{$10$}
\loigiai{
$9 + 3 = 12$.\\
\textbf{Đáp án của câu hỏi là A}
}
\end{ex}

\begin{ex} %[2D2V11]
Một chuồng có 14 con thỏ, người ta bán đi 5 con thỏ. Số thỏ còn lại trong chuồng là
\choice
{$8$ con}
{\True $9$ con}
{$10$ con}
{$7$ con}
\loigiai{
$14 - 5 = 9$ con.\\
\textbf{Đáp án của câu hỏi là B}
}
\end{ex}

\begin{ex} %[2D2C11]
Tìm số thích hợp điền vào ô trống: $18 - \dots = 9$.
\choice
{$8$}
{$7$}
{\True $9$}
{$10$}
\loigiai{
$18 - 9 = 9$.\\
\textbf{Đáp án của câu hỏi là C}
}
\end{ex}

\textbf{Phần 2. Trắc nghiệm đúng sai}
\setcounterpageref{ex}{0}

\begin{ex} %[2D2H11]
Xét tính đúng sai của các phép trừ:
\choiceTF
{\True $12 - 7 = 5$.}
{\True $15 - 8 = 7$.}
{$16 - 9 = 8$.}
{\True $16 - 9 = 7$.}
\loigiai{
\begin{itemize}
    \item $12 - 7 = 5$: Đúng.
    \item $15 - 8 = 7$: Đúng.
    \item $16 - 9 = 7 \neq 8$. Ý 3 Sai, Ý 4 Đúng.
\end{itemize}
\begin{itemize}
    \item a) \textbf{Đúng}.
    \item b) \textbf{Đúng}.
    \item c) \textbf{Sai}.
    \item d) \textbf{Đúng}.
\end{itemize}
}
\end{ex}

\begin{ex} %[2D2V11]
Cho biểu thức $M = 13 - 5$ và $N = 16 - 8$. Xét tính đúng sai:
\choiceTF
{\True $M = 8$.}
{\True $N = 8$.}
{$M > N$.}
{\True $M = N = 8$.}
\loigiai{
\begin{itemize}
    \item $M = 13 - 5 = 8$: Đúng.
    \item $N = 16 - 8 = 8$: Đúng.
    \item $M = N$ (không lớn hơn). Ý 3 Sai, Ý 4 Đúng.
\end{itemize}
\begin{itemize}
    \item a) \textbf{Đúng}.
    \item b) \textbf{Đúng}.
    \item c) \textbf{Sai}.
    \item d) \textbf{Đúng}.
\end{itemize}
}
\end{ex}

\textbf{Phần 3. Trả lời ngắn}
\setcounterpageref{ex}{0}

\begin{ex} %[2D2H11]
Tính: $11 - 8$.
\shortans{3}
\loigiai{
\textbf{Đáp số:} $3$\\
$11 - 8 = 3$.
}
\end{ex}

\begin{ex} %[2D2H11]
Tính: $14 - 6$.
\shortans{8}
\loigiai{
\textbf{Đáp số:} $8$\\
$14 - 6 = 8$.
}
\end{ex}

\begin{ex} %[2D2V11]
Tính: $17 - 9$.
\shortans{8}
\loigiai{
\textbf{Đáp số:} $8$\\
$17 - 9 = 8$.
}
\end{ex}

\begin{ex} %[2D2V11]
Tìm $x$ biết: $x + 8 = 14$.
\shortans{6}
\loigiai{
\textbf{Đáp số:} $6$\\
$x = 14 - 8 = 6$.
}
\end{ex}

\begin{ex} %[2D2V11]
Có 15 bông hoa, đem cắm vào bình 6 bông hoa. Hỏi còn lại bao nhiêu bông hoa chưa cắm?
\shortans{9}
\loigiai{
\textbf{Đáp số:} $9$\\
$15 - 6 = 9$ bông hoa.
}
\end{ex}

\begin{ex} %[2D2C11]
Tìm số thích hợp điền vào dấu ?: $16 - ? = 8$.
\shortans{8}
\loigiai{
\textbf{Đáp số:} $8$\\
$16 - 8 = 8$.
}
\end{ex}
